\section{Challenges and Future Directions}
\label{sec:outlook}
\subsection{Balancing Fidelity, Perceptual Quality, and Temporal Consistency}
Generative priors are most useful when observations are severely degraded, yet this is also where the reconstructed detail is most ambiguous. Strong conditioning can preserve scene structure while limiting texture synthesis; a more expressive sampling process can improve appearance but introduce spatial or temporal variation. The remaining challenge is to control which information comes from the observation and which is supplied by the prior, without reducing all three quality objectives to one fixed loss weight.

An effective approach should allocate uncertainty over coherent objects and trajectories. For example, several plausible reconstructions of an occluded surface can be acceptable, whereas its texture should remain consistent when it reappears within one reconstructed video. Trajectory-aware uncertainty, confidence-conditioned guidance, and explicit identity or structural constraints offer promising directions. Controllable restoration strength could further support different applications, allowing users to adjust perceptual detail while preserving selected attributes such as motion, text, or low-frequency structure.

\subsection{Reliable Temporal Modeling for Long Videos}
Image priors require mechanisms to connect independently modeled frames, while video priors bring motion knowledge learned from large datasets~\cite{svd,wan}. Both approaches face a distinction between plausible motion and the actual motion in the input. Correspondence becomes unreliable under heavy blur, darkness, occlusion, and atmospheric distortion, and attention can associate visually similar but unrelated regions. Combining correspondence uncertainty with learned temporal representations may improve robustness under these conditions.

Long videos create additional challenges. A model must maintain identity across temporary occlusion, recover context after viewpoint changes, and reset state at scene boundaries. Recurrent memory and attention caches reduce computation but can propagate early reconstruction errors. Future work should investigate selective state updates, motion-aware memory, and re-anchoring to current observations. Evaluation on long trajectories and scene transitions will be necessary to connect improvements in short-clip consistency to sustained performance.

\subsection{Unified and Fair Evaluation}
The diversity of backbones, pretraining data, degradation models, and inference settings makes it difficult to attribute improvements to a particular restoration technique. Controlled comparisons using shared initialization and matched data can reveal the value of conditioning or temporal modules, while complete-system comparisons capture practical advances from larger priors and optimized implementations. A common benchmark should support both forms of comparison and cover synthetic, real, specialized, and generated-video inputs.

Metric reliability remains equally important. Subjective assessment of diffusion-based VSR reveals different behavior across reference-based and no-reference predictors, and between within-source and overall evaluation~\cite{metric}. As generators produce sharper and smoother outputs, their remaining errors may involve identity, geometry, and motion rather than conventional blur or noise. New quality models should separate perceptual preference from source fidelity and temporal correctness, with evaluation across unseen restorers and degradations. Human judgments remain valuable for understanding these dimensions and calibrating objective measures.

\subsection{Broader Tasks and Physically Grounded Restoration}
Natural-video super-resolution has provided a productive setting for developing diffusion restoration. Weather removal, underwater enhancement, infrared restoration, face restoration, and event-guided processing demonstrate the broader potential of generative priors~\cite{difftta,weatherweaver,undive,hatir,flair,evdiffuser}. Extending these methods to diverse acquisition conditions requires a closer relationship between model design and image formation. Unknown mixtures such as scattering with motion blur, changing underwater illumination, or low-light compression can violate the assumptions of models trained on simpler degradations.

Physical modeling and multimodal observations provide complementary information. Events can constrain fast motion, reference images can preserve identity, and sensor metadata can inform noise or exposure models. Their integration should account for calibration errors, missing modalities, and changing conditions. HDR and infrared restoration also require objectives that preserve radiometric or thermal structure rather than merely reproducing the appearance of natural RGB video. Progress in these domains will depend on suitable data, models, and task-specific evaluation.

\subsection{Efficient Inference Beyond Sampling Steps}
One-step and few-step methods reduce the cost of repeated denoising, but high-resolution video processing can remain dominated by autoencoders, attention, tiling, and temporal overlap~\cite{dove,seedvr2,flash}. Efficiency therefore requires joint optimization of the full restoration system. Restoration-aware compression, lightweight temporal decoders, sparse attention, and reusable state offer opportunities to reduce cost while retaining detail and temporal coherence.

Streaming introduces a further balance between context and delay. Future frames improve reconstruction in offline processing, whereas online applications must operate with limited look-ahead. Adaptive computation could allocate more resources to difficult motion, occlusion, or degradation transitions while processing easier segments cheaply. Such strategies should be evaluated using first-output latency, sustained throughput, peak memory, and long-horizon quality, rather than denoising step count alone.

\subsection{Generalist and Adaptive Restoration Systems}
A single restoration model must accommodate degradations with different physical causes and potentially conflicting objectives. Prompts, expert routing, and shared representations provide ways to combine specialization with parameter sharing, but generalization to unfamiliar mixtures remains challenging. Continual adaptation and self-supervised learning may help models respond to changing cameras and environments without requiring paired clean video for every condition.

Agentic restoration extends adaptation to the choice and ordering of operators. Its effectiveness depends on diagnosing degradation, selecting compatible models, and deciding when further enhancement would damage content. Quality feedback can support this process, but an evaluator that rewards oversharpening or invented detail can also reinforce errors. Future systems should combine perceptual feedback with source-consistency constraints, maintain temporal decisions across a sequence, and avoid unnecessary repeated processing. For identity-sensitive and scientific applications, the ability to preserve trustworthy content is as important as the ability to generate convincing detail.

\section{Conclusion}
Diffusion models have expanded video restoration and enhancement through powerful image and video priors, flexible conditioning, and increasingly efficient inference. This survey has organized their development by prior source, learning strategy, temporal mechanism, task, and system paradigm. Across these categories, the central challenge is to connect synthesized detail to the observed scene while preserving coherent appearance and motion over time.

Further progress will require joint advances in restoration control, long-video modeling, domain-specific learning, and evaluation. Shared benchmarks can clarify the contribution of new methods, while calibrated metrics and human assessment can distinguish perceptual improvement from faithful reconstruction. Addressing these issues will support diffusion-based restoration systems that are both visually effective and reliable across diverse video applications.
